% =====================================================================
% Paper 44 (standalone): Operator-system extension of the BBB Krein
% spectral triple at finite cutoff; propagation number and envelope
% dependence on the Camporesi-Higuchi spinor bundle.
%
% Math.OA / J. Geom. Phys. style; sixth math.OA standalone in the
% GeoVac series. Siblings: Paper 38 (SU(2) propinquity convergence),
% Paper 39 (tensor-product propinquity convergence), Paper 40
% (universal 4/pi rate constant on compact Lie groups), Paper 42
% (Tomita-Takesaki / four-witness Riemannian closure), Paper 43
% (Lorentzian extension of the four-witness theorem).
%
% Status: first draft, May 2026, post-Sprint-L3a-1 (2026-05-17). Built
% from internal sprint memos:
%   debug/l3a_1_lorentzian_operator_system_memo.md
%   debug/l3_scoping_memo.md
%   debug/l3_literature_audit_memo.md
%   debug/sprint_l2_synthesis_memo.md
% =====================================================================
\documentclass[11pt,a4paper]{article}

\usepackage[T1]{fontenc}
\usepackage[utf8]{inputenc}
\usepackage{amsmath, amssymb, amsthm, mathrsfs, mathtools}
\usepackage{geometry}
\geometry{margin=1in}
\usepackage{hyperref}
\hypersetup{colorlinks=true, linkcolor=blue!50!black, citecolor=blue!50!black, urlcolor=blue!50!black}
% \usepackage{microtype}  % disabled: MiKTeX font-expansion environmental issue
\usepackage{graphicx}
\usepackage{booktabs}
\usepackage[numbers,sort&compress]{natbib}
\usepackage{xcolor}

\theoremstyle{plain}
\newtheorem{theorem}{Theorem}[section]
\newtheorem{lemma}[theorem]{Lemma}
\newtheorem{proposition}[theorem]{Proposition}
\newtheorem{corollary}[theorem]{Corollary}
\theoremstyle{definition}
\newtheorem{definition}[theorem]{Definition}
\newtheorem{remark}[theorem]{Remark}
\newtheorem{example}[theorem]{Example}

\newcommand{\nmax}{n_{\max}}
\newcommand{\Nt}{N_t}
\newcommand{\sthree}{S^{3}}
\newcommand{\Hilb}{\mathcal{H}}
\newcommand{\Krein}{\mathcal{K}}
\newcommand{\Op}{\mathcal{O}}
\newcommand{\Acal}{\mathcal{A}}
\newcommand{\Tcal}{\mathcal{T}}
\newcommand{\Bcal}{\mathcal{B}}
\newcommand{\opnorm}[1]{\lVert #1 \rVert_{\mathrm{op}}}
\newcommand{\Frob}[1]{\lVert #1 \rVert_{F}}
\newcommand{\norm}[1]{\lVert #1 \rVert}
\newcommand{\R}{\mathbb{R}}
\newcommand{\C}{\mathbb{C}}
\newcommand{\T}{\mathbb{T}}
\newcommand{\N}{\mathbb{N}}
\newcommand{\Z}{\mathbb{Z}}
\newcommand{\Q}{\mathbb{Q}}
\newcommand{\SU}{\mathrm{SU}}
\newcommand{\SO}{\mathrm{SO}}
\newcommand{\Vol}{\mathrm{Vol}}
\newcommand{\Tr}{\mathrm{Tr}}
\newcommand{\Mat}{M}
\newcommand{\DCH}{D_{\mathrm{CH}}}
\newcommand{\DGV}{D_{\mathrm{GV}}}
\newcommand{\DL}{D_L}
\newcommand{\JL}{J_L}
\newcommand{\HGV}{\mathcal{H}_{\mathrm{GV}}}
\newcommand{\Cl}{\mathrm{Cl}}
\newcommand{\prop}{\mathrm{prop}}
\newcommand{\TT}{\mathrm{TT}}
\newcommand{\KMS}{\mathrm{KMS}}
\newcommand{\BW}{\mathrm{BW}}
\newcommand{\spat}{\mathrm{spat}}
\newcommand{\temp}{\mathrm{temp}}
\newcommand{\Weyl}{\mathrm{Weyl}}
\newcommand{\achievable}{\mathrm{ach}}

\title{Operator-system extension of the BBB Krein spectral triple\\
at finite cutoff:\ propagation number and\\
envelope dependence on the Camporesi--Higuchi spinor bundle}
\author{J.~Loutey\thanks{Independent researcher. \texttt{jloutey@gmail.com}.
This work was developed in an AI-augmented agentic workflow with
Anthropic's Claude. Implementation, internal memos, and bibliographic
collation were performed collaboratively; all theorems, design choices,
and editorial decisions are the author's responsibility.}}
\date{May 2026 (preprint)}

\begin{document}
\maketitle

\begin{abstract}
The Connes--van~Suijlekom~\cite{connes_vs2021} spectral truncation of
$C^\infty(\sthree)$ on the round $\sthree = \SU(2)$ at cutoff $\nmax$
produces a $*$-closed but \emph{not} multiplicatively closed linear
subspace $\Op_{\nmax} \subset \Mat_{N}(\C)$ with propagation number
$\prop(\Op_{\nmax}) = 2$ (Paper~32~\S~III, matching the Connes--vS
worked example on the Toeplitz operator system $C(S^1)^{(n)}$
verbatim). We construct the natural Lorentzian extension
$\Op_{\nmax, \Nt}^L \subset \Bcal(\Krein_{\nmax, \Nt})$ on the
Bizi--Brouder--Besnard~\cite{bizi_brouder_besnard2018} Krein
spectral triple at signature $(s, t) = (3, 1)$, $(m, n) = (4, 6)$,
in the Peskin--Schroeder chiral basis with Lorentzian Dirac
$\DL = i\bigl(\gamma^0 \otimes \partial_t + \DGV \otimes I_{\Nt}\bigr)$
per van~den~Dungen~2016~\cite{vandungen2016} Proposition~4.1, and the
Camporesi--Higuchi~\cite{camporesi_higuchi1996} spatial spinor bundle.
Multipliers are tensor products of chirality-doubled spatial scalar
$3$-$Y$ multipliers with a polynomial-basis temporal subalgebra on a
finite uniform $t$-grid. At $\Nt = 1$ the construction reduces
bit-identically to the Riemannian chirality-doubled operator system on
the FullDirac Hilbert space, with Frobenius residual exactly $0.0$ in
float64. We prove $\Op_{\nmax, \Nt}^L$ is a Lorentzian truncated
operator system (it is $*$-closed, contains the identity, and is not
multiplicatively closed) and exhibit a witness pair lifting the
Paper~32~\S~III witness verbatim. The principal structural finding is
that the propagation number is \emph{envelope-dependent}:\ targeting
the natural Weyl-doubled achievable envelope $\dim_\Weyl^2 \cdot \Nt$,
we obtain $\prop(\Op_{\nmax, \Nt}^L) = 2$ for every $\nmax \ge 2$,
matching Paper~32~\S~III verbatim; targeting the full $\Bcal(\Krein)$
envelope $\dim_\Krein^2$, we obtain $\prop = \infty$ at every cell.
The Weyl-doubling and the commutative temporal subalgebra
\emph{block} scalar multipliers from generating chirality-flipping or
non-diagonal temporal operators; the achievable envelope is the
correct natural target for scalar-multiplier closure on the spinor
bundle. A separate substantive finding is that the Krein-positive
restriction is trivial:\ every chirality-doubled scalar multiplier
$M \oplus M$ commutes with the chirality-swap fundamental symmetry
$\JL$, so $\Op^{L,+} = \Op^L$ exactly. The Krein-positive program is
shifted from the operator-system substrate (where it is automatic) to
the state level, which is the natural setting for any future
Lorentzian-propinquity convergence argument. The construction is the
foundation on which any subsequent continuum Lorentzian-propinquity
program (Sprint L3b, see~\S~\ref{sec:open}) would build. Honest scope:
this paper closes the operator-system substrate at finite cutoff;
continuum / GH-limit / Lorentzian-propinquity convergence is open. No
published Lorentzian propinquity construction exists as of May 2026:\
Latr\'emoli\`ere~2017/2023~\cite{latremoliere_metric_st_2017} and
Hekkelman--McDonald~2024~\cite{hekkelman_mcdonald2024} are strictly
Riemannian.
\end{abstract}

\bigskip\noindent
\textbf{MSC2020 classifications:}\ 58B34 (noncommutative geometry),
46L07 (operator systems), 46C20 (spaces with indefinite inner product),
46L60 (applications of selfadjoint operator algebras to physics),
46L87 (noncommutative Riemannian and metric geometry).

\medskip\noindent
\textbf{Keywords:}\ Connes--van~Suijlekom truncated operator system,
propagation number, Krein spectral triple, Bizi--Brouder--Besnard
indefinite signature, Camporesi--Higuchi Dirac, Lorentzian
noncommutative geometry, spectral truncation, finite-cutoff
operator-system substrate.

% =====================================================================
\section{Introduction}\label{sec:intro}
% =====================================================================

A central technical object in the spectral-truncation framework of
Connes and van~Suijlekom~\cite{connes_vs2021} is the truncated operator
system
\begin{equation}\label{eq:cvs_riem}
   \Op_\Lambda \;=\; P_\Lambda\, C^\infty(M)\, P_\Lambda
   \;\subset\; \Bcal(\Hilb_\Lambda),
\end{equation}
where $P_\Lambda$ projects $L^2(M)$ onto the span of the first
$\Lambda$ eigenfunctions of the Dirac operator (or, equivalently in
the GeoVac convention, onto the first $\nmax$ Fock shells of the
Fock-projected spectrum). On the round $\sthree = \SU(2)$ with
spinor-harmonic basis adapted to the SO(4) decomposition, this
construction was carried out in Paper~32~\S~III of the present
author's framework documentation and verified to have propagation
number
\begin{equation}\label{eq:prop_riem}
   \prop(\Op_{\nmax}) \;=\; 2 \qquad (\nmax \in \{2, 3, 4\}),
\end{equation}
matching the Connes--vS worked example on the Toeplitz operator
system $C(S^1)^{(n)}$ verbatim~\cite{connes_vs2021}. This was the
strongest single quantitative alignment between GeoVac and the
published Connes--vS framework as of the Round-2 sprint
(\texttt{debug/wh1\_round2\_propagation\_memo.md}, 2026-05-03).

The companion Paper~43~\cite{paper43} of the present author closes
the four-witness Wick-rotation theorem
(Hartle--Hawking~\cite{hartle_hawking1976},
Sewell~\cite{sewell1982},
Bisognano--Wichmann~\cite{bisognano_wichmann1976},
Unruh~\cite{unruh1976}) at the operator-system level on the
Lorentzian Krein spectral triple at signature $(3, 1)$ in the
Bizi--Brouder--Besnard~\cite{bizi_brouder_besnard2018} classification
at $(m, n) = (4, 6)$, with the modular Hamiltonian and Tomita
structure on a hemispheric wedge. Paper~43 takes the operator-system
substrate as input;\ the present paper constructs it.

\subsection{What this paper does}

We define and study the natural Lorentzian extension
\begin{equation}\label{eq:cvs_lor}
   \Op_{\nmax, \Nt}^L
   \;=\; P_{\nmax, \Nt} \cdot \bigl[\,
                C^\infty(\sthree) \otimes C^\infty(\R_t)_{\mathrm{cutoff}}
              \,\bigr]
   \cdot P_{\nmax, \Nt}
   \;\subset\; \Bcal(\Krein_{\nmax, \Nt})
\end{equation}
on the Krein space
$\Krein_{\nmax, \Nt} = \HGV^{\nmax} \otimes \C^{\Nt}$ of Paper~43
\S~\ref{sec:krein_space} (Sprint L2-B). Here $\HGV^{\nmax}$ is the
chirality-doubled Camporesi--Higuchi~\cite{camporesi_higuchi1996}
spinor space at Fock cutoff $\nmax$, with
$\dim_\C \HGV^{\nmax} = \tfrac{2}{3}\nmax(\nmax + 1)(\nmax + 2)
\in \{4, 16, 40, 80, 140\}$ at $\nmax \in \{1, 2, 3, 4, 5\}$, and
$\C^{\Nt}$ is the temporal $L^2$ cutoff on the uniform grid
$t_k = -T_{\max} + k \cdot 2T_{\max}/(\Nt - 1)$. Spatial scalar
multipliers act on $\HGV^{\nmax}$ via the chirality-block-doubled
spinor lift (Clebsch--Gordan decomposition $\oplus$ Avery--Wen--Avery
$3$-$Y$ integrals); temporal multipliers act on $\C^{\Nt}$ by
pointwise multiplication via the polynomial basis $g_p(t) = t^p$ for
$p = 0, \ldots, \Nt - 1$.

The principal results are:

\begin{enumerate}\setlength{\itemsep}{2pt}
\item \textbf{Operator-system axioms hold} (Proposition~\ref{prop:axioms},
Section~\ref{sec:op_system}). $\Op_{\nmax, \Nt}^L$ contains the
identity, is closed under conjugate transposition, and is \emph{not}
multiplicatively closed:\ an explicit witness pair lifting the
Paper~32 spatial witness produces a product $ab$ with Frobenius
residual $\sim 38\%$ relative to the linear span of $\Op^L$
(Theorem~\ref{thm:witness_pair}, Section~\ref{sec:witness}).
\item \textbf{Riemannian limit at $\Nt = 1$ is bit-exact}
(Theorem~\ref{thm:riemannian_limit}, Section~\ref{sec:riem_limit}). The
construction reduces to the chirality-doubled FullDirac operator
system on $\HGV^{\nmax}$ with maximum residual $0.0$ in float64 at
every $\nmax \in \{1, 2, 3\}$, verifying compatibility with the
Sprint L2-B Krein-space basis convention.
\item \textbf{Envelope-dependent propagation number} (Theorem~%
\ref{thm:prop_envelope_dependent}, Section~\ref{sec:propagation}).
Targeting the Weyl-doubled \emph{achievable} envelope
$\dim_\achievable = \dim_\Weyl^2 \cdot \Nt$,
\begin{equation}\label{eq:prop_lor_ach_intro}
   \prop_{\achievable}(\Op_{\nmax, \Nt}^L) \;=\; 2
   \qquad \text{for every } \nmax \ge 2, \; \Nt \ge 1,
\end{equation}
verbatim with Paper~32~\S~III. Targeting the full $\Bcal(\Krein)$
envelope $\dim_\Krein^2$,
\begin{equation}\label{eq:prop_lor_full_intro}
   \prop_{\mathrm{full}}(\Op_{\nmax, \Nt}^L) \;=\; \infty
   \qquad \text{at every cell.}
\end{equation}
The discrepancy is a real structural feature of the chirality-doubled
spinor bundle:\ scalar multipliers $M \oplus M$ cannot generate
chirality-flipping operators (off-diagonal blocks of the chirality
splitting), and a commutative temporal subalgebra cannot generate
non-diagonal temporal operators. Eq.~\eqref{eq:prop_lor_ach_intro} is
the natural Connes--vS-style propagation question for scalar
multipliers on a chirality-doubled bundle;\
\eqref{eq:prop_lor_full_intro} is a strictly larger naive envelope not
reached by the scalar-multiplier construction.
\item \textbf{Krein-positive restriction is trivial}
(Proposition~\ref{prop:krein_positivity_trivial},
Section~\ref{sec:krein_positivity}). Every multiplier in
$\Op_{\nmax, \Nt}^L$ commutes with the chirality-swap fundamental
symmetry $\JL = \gamma^0 \otimes I_{\Nt}$, so it preserves the
Krein-positive cone $\Krein^+ = \mathrm{ker}(\JL - I)$. The natural
restriction $\Op^{L,+}$ to Krein-positive preservers equals
$\Op^L$ itself.
\end{enumerate}

\subsection{Position in the literature}

The Lorentzian-propinquity question
sits at the intersection of three threads with empty intersection as
of May 2026:\
\begin{enumerate}\setlength{\itemsep}{2pt}
\item[(a)] \textbf{Riemannian propinquity convergence on spectral
truncations.} Connes--van~Suijlekom~\cite{connes_vs2021} introduces
the truncated operator system and the propagation-number invariant;
Latr\'emoli\`ere~\cite{latremoliere_metric_st_2017,latremoliere2018}
develops the propinquity for metric spectral triples;\ the
explicit Gromov--Hausdorff convergence on a non-trivial example was
deferred to ``elsewhere'' three times in \cite{connes_vs2021}
(lines~220--221, 1101, 2616--2617 of arXiv:2004.14115v2). Subsequent
work by
Hekkelman~\cite{hekkelman2022},
Hekkelman--McDonald~\cite{hekkelman_mcdonald2024,
hekkelman_mcdonald2024b},
Leimbach--van~Suijlekom~\cite{leimbach_vs2024},
Toyota~\cite{toyota2023},
Farsi--Latr\'emoli\`ere~\cite{farsi_latremoliere2024,
farsi_latremoliere2025} addresses the flat structures $S^1$, $T^d$,
and Berezin--Toeplitz $S^2$. The first non-abelian case
$\sthree = \SU(2)$ was closed by the present author in Paper~38
\cite{paper38}; the tensor-product extension was closed in
Paper~39~\cite{paper39}; the universal compact-Lie-group rate constant
$4/\pi$ was closed in Paper~40~\cite{paper40_unified}. All of these
are strictly Riemannian.
\item[(b)] \textbf{Lorentzian noncommutative geometry.}
Strohmaier~\cite{strohmaier2006},
van~den~Dungen~\cite{vandungen2016},
Franco--Eckstein~\cite{franco_eckstein2014},
Bizi--Brouder--Besnard~\cite{bizi_brouder_besnard2018},
Devastato--Lizzi--Martinetti~\cite{devastato_lizzi_martinetti2018},
and
Nieuviarts~\cite{nieuviarts2025a} construct indefinite (Krein)
spectral triples and address Wick rotation at the level of the data
$(\Acal, \Krein, \DL, \JL)$. None of these works carries the
construction through to a Connes--van~Suijlekom truncated operator
system on the Krein side, nor to a propagation-number invariant.
\item[(c)] \textbf{Lorentzian Gromov--Hausdorff convergence.} The
recent Mondino--Sämann program on synthetic Lorentzian
GH~\cite{mondino_samann2025} and the
Bykov--Minguzzi--Suhr work on Lorentzian convergence of metric
spaces~\cite{bykov_minguzzi_suhr2024} develop a metric-space-side
notion of Lorentzian convergence but do not interface with the
operator-system/spectral-triple side.
\end{enumerate}
The present paper closes the operator-system \emph{substrate} of any
prospective Lorentzian-propinquity convergence program at finite
cutoff;\ it does not close the propinquity itself. The substrate is the
foundation on which a future propinquity argument would be built.

\subsection{Why this is a separate paper from Paper~43}

Paper~43~\cite{paper43} closes the four-witness Wick-rotation theorem
at the \emph{modular-Hamiltonian} level on a hemispheric wedge of the
Lorentzian Krein space, using the same Krein construction
$\Krein_{\nmax, \Nt}$, Lorentzian Dirac $\DL$, and BBB axiom audit at
$(m, n) = (4, 6)$ as the present paper. The Paper~43 results
(integer spectrum of the geometric BW-$\alpha$ generator, $\beta$-%
independent wedge KMS state, Tomita--Takesaki polar decomposition,
flow conjugacy, six-witness collapse) operate at the level of
\emph{specific operators} on the Krein space:\ the wedge projector,
the modular operator $\Delta$, the modular Hamiltonian $K_L^\alpha$
and $K_L^{\TT}$. They do not directly study the linear span of all
spatial-times-temporal multipliers as a $*$-closed subspace of
$\Bcal(\Krein)$.

The present paper studies that subspace explicitly. The
propagation-number invariant is independent of any wedge structure or
modular flow:\ it is a statement about the multiplicative closure of
the operator system as a whole, and its envelope-dependence at the
Weyl-doubled vs full $\Bcal(\Krein)$ level is a separate structural
feature not visible in Paper~43's hemispheric-wedge analysis.
Conceptually, Paper~43 is the analog of computing modular data on a
specific algebra (here the wedge-restricted operator system); the
present paper is the analog of constructing the algebra itself
together with its envelope-relative structural invariants.
Computationally, the two papers share their lower layers
(\texttt{geovac/krein\_space\_construction.py},
\texttt{geovac/lorentzian\_dirac.py},
\texttt{geovac/connes\_axiom\_audit\_31.py}) but Paper~43 does not
import \texttt{geovac/operator\_system\_lorentzian.py} (the Sprint
L3a-1 module) and vice versa.

\subsection{Outline}

Section~\ref{sec:setup} recalls the Krein-space construction at
$(m, n) = (4, 6)$ from Paper~43 and fixes notation.
Section~\ref{sec:op_system} defines the Lorentzian truncated operator
system $\Op^L$, proves $*$-closure and identity-membership, and
discusses the multiplier bookkeeping.
Section~\ref{sec:riem_limit} verifies the Riemannian-limit recovery at
$\Nt = 1$ (load-bearing falsifier).
Section~\ref{sec:propagation} proves the envelope-dependent
propagation theorem.
Section~\ref{sec:krein_positivity} establishes the trivial
Krein-positive restriction.
Section~\ref{sec:witness} exhibits the witness pair.
Section~\ref{sec:wedge} connects to the Paper~43 hemispheric wedge.
Section~\ref{sec:open} states open questions and the recommended
first Sprint L3b move.
Section~\ref{sec:conclusion} concludes.

% =====================================================================
\section{Setup: Krein space, Lorentzian Dirac, BBB axiom audit}\label{sec:krein_space}
\label{sec:setup}
% =====================================================================

Sections~\ref{sec:krein}--\ref{sec:bbb} recall the data of the
Lorentzian Krein spectral triple from
Paper~43~\cite{paper43}~\S~3 (Sprints L2-B/C/D). The reader familiar
with Paper~43 may skim or skip these subsections.

\subsection{The Krein space $\Krein_{\nmax, \Nt}$}\label{sec:krein}

Fix a Fock cutoff $\nmax \in \N$ and a temporal grid size $\Nt \in
\N$ on a half-width $T_{\max} > 0$. The Krein space at cutoff is
\begin{equation}\label{eq:krein_def}
   \Krein_{\nmax, \Nt}
   \;:=\; \HGV^{\nmax} \otimes \C^{\Nt},
\end{equation}
where $\HGV^{\nmax}$ is the chirality-doubled Camporesi--Higuchi
spinor space at cutoff $\nmax$
(Paper~32~\S~III~\cite{paper32}, Paper~43~\S~3.1) with basis labelled
$\bigl(n_{\mathrm{fock}}, l, 2 m_j, \mathrm{chirality}\bigr)$ in the
\texttt{FullDiracLabel} convention. The spatial dimension is
\begin{equation}\label{eq:dim_spatial}
   \dim_\C \HGV^{\nmax}
   \;=\; \tfrac{2}{3}\,\nmax(\nmax + 1)(\nmax + 2)
   \;=\; \begin{cases} 4 & \nmax = 1 \\ 16 & \nmax = 2 \\
         40 & \nmax = 3 \\ 80 & \nmax = 4 \\ 140 & \nmax = 5 \end{cases}
\end{equation}
and the total Krein dimension is $\dim_\C \Krein_{\nmax, \Nt} =
\dim_\C \HGV^{\nmax} \cdot \Nt$.

We use the symmetric temporal grid
\begin{equation}\label{eq:t_grid}
   t_k \;:=\; -T_{\max} + k \cdot \tfrac{2 T_{\max}}{\Nt - 1}
   \qquad (k = 0, \ldots, \Nt - 1),
\end{equation}
with the convention $t_0 = 0$ at $\Nt = 1$. The fundamental symmetry
on $\Krein_{\nmax, \Nt}$ is
\begin{equation}\label{eq:J_def}
   \JL \;:=\; \gamma^0 \otimes I_{\Nt},
\end{equation}
where $\gamma^0$ acts as chirality-swap on $\HGV$ in the
Peskin--Schroeder chiral basis (West-coast metric $\eta =
\mathrm{diag}(+, -, -, -)$). The basic Krein-space axioms
$\JL^2 = +I$, $\JL^\dagger \JL = I$, $\JL^\dagger = \JL$, and the
splitting $\Krein = \Krein^+ \oplus \Krein^-$ into $\pm 1$-eigenspaces
of $\JL$ are verified bit-exactly in Sprint L2-B (Paper~43~\S~3.1).

\subsection{The Lorentzian Dirac $\DL$}\label{sec:lor_dirac}

The Lorentzian Dirac operator on $\Krein_{\nmax, \Nt}$ is
\begin{equation}\label{eq:DL_def}
   \DL \;:=\; i\bigl(
              \gamma^0 \otimes \partial_t
              \;+\; \DGV \otimes I_{\Nt}
            \bigr),
\end{equation}
per van~den~Dungen~2016~\cite{vandungen2016} Proposition~4.1 with
$(s, t) = (3, 1)$ giving $i^t = i^1 = i$, where $\DGV$ is the
truthful Camporesi--Higuchi Dirac on $\HGV^{\nmax}$
(spectrum $|\lambda_n| = n + \tfrac{1}{2}$ with multiplicities
$g_n^{\mathrm{Dirac}} = 2(n+1)(n+2)$ on each chirality) and
$\partial_t$ is the centered finite-difference operator on the
uniform grid \eqref{eq:t_grid} with Dirichlet zero boundary
conditions. The sign $i^t = +i$ is structurally forced by
Krein-self-adjointness $\DL^\times = \DL$ (Paper~43~\S~3.2,
\texttt{geovac/lorentzian\_dirac.py}). At $\Nt = 1$, $\partial_t = 0$
and $\DL = i\DGV$ reduces to the Riemannian Camporesi--Higuchi
Dirac up to a global imaginary scaling that drops out of all
load-bearing falsifiers (Paper~43~\S~3.2).

\subsection{The BBB axiom audit at $(m, n) = (4, 6)$}\label{sec:bbb}

The Bizi--Brouder--Besnard~\cite{bizi_brouder_besnard2018}
classification of indefinite spectral triples by
$(m, n) = (t + s, t - s)$ mod $8$ places Lorentzian $(3, 1)$ West-coast
at $(m, n) = (4, 6)$, with primitive signs (BBB Table~1)
\begin{equation}\label{eq:bbb_signs}
   \varepsilon = +1, \quad
   \varepsilon'' = -1, \quad
   \kappa = -1, \quad
   \kappa'' = +1.
\end{equation}
The lift of the Cl$(3,1)$ charge-conjugation matrix $U_4 = i\gamma^2$
to $\Krein$ is $\JL = J_{L, \mathrm{spatial}} \otimes K_t$ with $K_t$
the identity-unitary times complex conjugation on $\C^{\Nt}$
(Paper~43~\S~3.3, \texttt{geovac/connes\_axiom\_audit\_31.py}). The
four BBB-predicted-sign axioms
\begin{equation}\label{eq:bbb_axioms}
   \JL^2 = +I, \quad
   \{\JL, \gamma^5\} = 0, \quad
   \{\JL, \eta\} = 0, \quad
   \JL \DL = +\DL \JL
\end{equation}
pass bit-exactly (Frobenius residual $= 0.0$ in float64) at every
tested $(\nmax, \Nt) \in \{1, 2, 3\} \times \{1, 11, 21\}$
(Paper~43~\S~3.3). The non-load-bearing universal BBB axiom
$\{\gamma^5, \DL\} = 0$ \emph{fails} on truthful $\DGV$ for the
reason described in Paper~43~\S~9.3:\ this is a known structural
feature (chirality-diagonal Camporesi--Higuchi Dirac in the
Peskin--Schroeder chiral basis), not a bug, and is preserved by the
present construction without affecting any of its load-bearing
results.

\subsection{Spatial scalar multipliers on the chirality-doubled bundle}
\label{sec:spat_mult}

The Riemannian Connes--vS multiplier $M_{NLM}$ on $L^2(\sthree)$
acts by pointwise multiplication of a scalar hyperspherical harmonic
$Y^{(3)}_{NLM}$:\
\begin{equation}\label{eq:M_NLM_riem}
   \langle (n,l,m_l) \mid M_{NLM} \mid (n', l', m_l') \rangle
   \;=\; \int_{\sthree}
           \overline{Y_{n l m_l}}\, Y^{(3)}_{NLM}\, Y_{n' l' m_l'}\, dV,
\end{equation}
with selection rules $|n - n'| + 1 \le N \le n + n' - 1$ (SO(4)
triangle), $|l - l'| \le L \le l + l'$ with $l + l' + L$ even
(parity), and $M = m_l - m_l'$. The integral is the
Avery--Wen--Avery~\cite{avery_wen_avery2002} 3-Y integral, evaluated
symbolically in exact $\Q$-arithmetic via
\texttt{geovac/so4\_three\_y\_integral.py}.

The lift to the Weyl-spinor sector $\HGV^{\Weyl} \subset \HGV$ is
the Clebsch--Gordan-projected matrix
\begin{equation}\label{eq:M_NLM_weyl}
   \langle (n, l, j, m_j) \mid M_{NLM} \mid
            (n', l', j', m_j') \rangle
   \;=\; \sum_{m_l, m_l'} C^{j, m_j}_{l, m_l; 1/2, m_s}
                     \, C^{j', m_j'}_{l', m_l'; 1/2, m_s}
                     \, \delta_{m_s, m_j - m_l}\,
                       (\text{Eq.}~\eqref{eq:M_NLM_riem}),
\end{equation}
via \texttt{geovac/spinor\_operator\_system.py}. Both chirality
sectors of $\HGV = \HGV^{\Weyl} \oplus \HGV^{\Weyl}$ see the
scalar multiplier identically, so the FullDirac multiplier is
\begin{equation}\label{eq:M_NLM_full}
   M^{\spat}_{NLM} \;:=\; M^{\Weyl}_{NLM} \oplus M^{\Weyl}_{NLM}
   \;\in\; \Bcal(\HGV^{\nmax}),
\end{equation}
i.e.\ block-diagonal in the two chirality sectors with equal blocks.
This is exactly the \texttt{FullDiracTruncatedOperatorSystem.build\_%
full\_dirac\_multiplier\_matrix} construction of
Sprint WH1-R3.5~\cite{paper32}.

\subsection{Temporal multipliers on $\C^{\Nt}$}\label{sec:temp_mult}

We choose the polynomial basis $g_p(t) = t^p$ for
$p = 0, 1, \ldots, \Nt - 1$, evaluated on the temporal grid
\eqref{eq:t_grid}. The corresponding diagonal matrix is
\begin{equation}\label{eq:M_temp}
   M_p^{\temp}
   \;:=\; \mathrm{diag}\bigl(g_p(t_0), g_p(t_1), \ldots,
                             g_p(t_{\Nt - 1})\bigr)
   \;\in\; \Bcal(\C^{\Nt}).
\end{equation}
The Vandermonde matrix
$V_{kp} = t_k^p$ is invertible at $\Nt$ distinct grid points, so the
$\Nt$-element family $\{M_p^{\temp}\}_{p = 0}^{\Nt - 1}$ spans the
full $\Nt$-dimensional commutative subalgebra of diagonal matrices on
$\C^{\Nt}$. At $\Nt = 1$, only the constant $g_0(t) = 1$ is
available, and $M_0^{\temp} = 1$ is the $1 \times 1$ identity.

% =====================================================================
\section{The Lorentzian truncated operator system}\label{sec:op_system}
% =====================================================================

\begin{definition}[Lorentzian truncated operator system]
\label{def:O_L}
At fixed $(\nmax, \Nt, T_{\max})$, the \emph{Lorentzian truncated
operator system} on the Krein space $\Krein_{\nmax, \Nt}$ is the
$\C$-linear span
\begin{equation}\label{eq:O_L_def}
   \Op_{\nmax, \Nt}^L
   \;:=\; \mathrm{span}_\C
         \Bigl\{
           M^{\spat}_{NLM} \otimes M^{\temp}_p
           \;:\;
           (N, L, M) \in \mathcal{S}_{\spat}(\nmax),
           \;\; 0 \le p \le \Nt - 1
         \Bigr\}
   \;\subset\; \Bcal(\Krein_{\nmax, \Nt}),
\end{equation}
where $\mathcal{S}_{\spat}(\nmax)$ is the set of multiplier labels
$(N, L, M)$ admitting a nonzero spinor-lift matrix
\eqref{eq:M_NLM_full} at cutoff $\nmax$ (SO(4) triangle, parity,
$M$-conservation, plus the $\nmax$-truncation kernel filter
$|n - n'| + 1 \le N \le n + n' - 1$ both ways).
\end{definition}

The number of spatial multipliers is exactly
$\dim_\C \Op_{\nmax} = 1, 14, 55, 140, 285$ at $\nmax = 1, 2, 3, 4, 5$
(Paper~32~\S~III~\cite{paper32}), and the total number of generators of
$\Op^L$ is $\dim_\C \Op_{\nmax} \cdot \Nt$.

\begin{proposition}[Operator-system axioms]\label{prop:axioms}
At every $(\nmax, \Nt)$ with $\Nt \ge 1$:
\begin{enumerate}\setlength{\itemsep}{2pt}
\item[(i)] $\Op_{\nmax, \Nt}^L$ is closed under conjugate
transposition $A \mapsto A^\dagger$ on $\Bcal(\Krein)$.
\item[(ii)] $I_{\Krein} \in \Op_{\nmax, \Nt}^L$, with $I_{\Krein}$
realised by the trivial multiplier $M^{\spat}_{0,0,0} \otimes
M_0^{\temp} = I_{\dim_\spat} \otimes I_{\Nt}$.
\item[(iii)] $\Op_{\nmax, \Nt}^L$ is \emph{not} multiplicatively
closed for any $\nmax \ge 2$ (witness pair, Theorem~%
\ref{thm:witness_pair}, Section~\ref{sec:witness}).
\end{enumerate}
\end{proposition}

\begin{proof}
For (i), the spatial conjugation
\eqref{eq:M_NLM_riem} satisfies the Wigner relation
\begin{equation}\label{eq:Y_conj}
   \overline{Y^{(3)}_{NLM}} \;=\; (-1)^{M_+}\, Y^{(3)}_{N, L, -M},
\end{equation}
for a definite sign $(-1)^{M_+}$ depending on the real-spherical
convention (we use complex sphericals so $M_+ = M$); the lifted
spinor matrix \eqref{eq:M_NLM_full} inherits this and
$(M^{\spat}_{NLM})^\dagger = \pm M^{\spat}_{N, L, -M}$. The temporal
multiplier $M_p^{\temp}$ is real-diagonal hence self-adjoint:
$(M_p^{\temp})^\dagger = M_p^{\temp}$. Tensor products of $*$-closed
families are $*$-closed:
\[
   (M^{\spat}_{NLM} \otimes M_p^{\temp})^\dagger
   \;=\; (M^{\spat}_{NLM})^\dagger \otimes M_p^{\temp}
   \;=\; \pm\, M^{\spat}_{N, L, -M} \otimes M_p^{\temp}
   \;\in\; \Op^L,
\]
provided $(N, L, -M) \in \mathcal{S}_{\spat}(\nmax)$, which follows
from the $M$-symmetry of the SO(4) selection rules.

For (ii), the trivial multiplier $M^{\spat}_{0,0,0}$ is the
$\HGV^{\nmax}$-identity (the constant function $Y^{(3)}_{0,0,0}$
acts as multiplication by a nonzero scalar; under the standard
normalisation it acts as the identity matrix), and
$M_0^{\temp}$ is the $\C^{\Nt}$-identity by
$g_0(t) \equiv 1$. The tensor $I_{\dim_\spat} \otimes I_{\Nt} =
I_{\Krein}$.

For (iii) we postpone to Section~\ref{sec:witness}.
\end{proof}

\begin{remark}[Operator-system terminology]\label{rem:op_system}
By standard usage~\cite{paulsen2002,connes_vs2021}, an
\emph{operator system} is a $*$-closed unital linear subspace of
$\Bcal(\Hilb)$ (without the requirement of multiplicative closure).
Proposition~\ref{prop:axioms} establishes that
$\Op_{\nmax, \Nt}^L$ is an operator system in this sense. The
absence of multiplicative closure (item (iii)) is the structural
feature that makes the operator-system substrate \emph{strictly
weaker} than the $C^*$-envelope $C^*(\Op^L) \subset \Bcal(\Krein)$;\
the operator-system data is precisely the algebraic content surviving
the spectral truncation that the $C^*$-envelope discards (cf.\
Connes--vS \cite{connes_vs2021}, ``if one takes the $C^*$-algebra
generated by $PAP$ one gets a non-informative full matrix algebra'').
\end{remark}

% =====================================================================
\section{Riemannian limit at $\Nt = 1$}\label{sec:riem_limit}
% =====================================================================

The load-bearing falsifier for the Lorentzian extension is that at
$\Nt = 1$ — where the temporal slot reduces to a single point and the
construction reduces to a static spatial slice — the operator-system
substrate must recover the chirality-doubled Riemannian operator
system bit-identically.

\begin{theorem}[Bit-exact Riemannian-limit recovery at $\Nt = 1$]
\label{thm:riemannian_limit}
At $\Nt = 1$, the Lorentzian truncated operator system reduces to the
chirality-doubled Riemannian operator system on $\HGV^{\nmax}$:
\begin{equation}\label{eq:riem_limit}
   \Op_{\nmax, 1}^L
   \;=\; \Op_{\nmax}^{\mathrm{Full\,Dirac}}
   \;=\; \bigl\{
           M \oplus M \;:\; M \in \Op_{\nmax}^{\Weyl}
         \bigr\}
   \;\subset\; \Bcal(\HGV^{\nmax}).
\end{equation}
The reduction is bit-exact in float64:\ for every multiplier label
$(N, L, M, p = 0) \in \mathcal{S}_{\spat}(\nmax) \times \{0\}$, the
Lorentzian multiplier matrix \eqref{eq:M_NLM_full} equals the
corresponding chirality-doubled spatial multiplier on $\HGV^{\nmax}$
with maximum Frobenius residual $= 0.0$ at every tested
$\nmax \in \{1, 2, 3\}$.
\end{theorem}

\begin{proof}[Proof and computational verification]
At $\Nt = 1$, $M_0^{\temp} = 1$ is the $1 \times 1$ identity scalar,
so the tensor product $M^{\spat}_{NLM} \otimes M_0^{\temp} =
M^{\spat}_{NLM}$ is bit-identical to the chirality-doubled spatial
multiplier on $\HGV^{\nmax}$. The full list of multiplier labels
$\mathcal{S}_{\spat}(\nmax)$ at $\Nt = 1$ matches the
FullDiracTruncatedOperatorSystem multiplier list one-for-one
(same SO(4) triangle, parity, and $M$-conservation filters), and the
matrix construction calls the same
\texttt{build\_spinor\_multiplier\_matrix} routine on the same
Avery--Wen--Avery 3-Y integrals.

The computational verification is documented in
\texttt{tests/test\_operator\_system\_lorentzian.py::%
test\_riemannian\_limit\_bit\_exact}: at $\nmax \in \{1, 2, 3\}$,

\medskip
\begin{center}\renewcommand{\arraystretch}{1.20}
\begin{tabular}{c|c|c|c|c}
$\nmax$ & $\dim_\Krein$ & $\#$ multipliers
   & $\dim(\Op^L)$ & $\max \Frob{M_{\mathrm{self}} - M_{\mathrm{ref}}}$ \\
\hline
$1$ & $4$ & $1$ & $1$ & $0.0$ (exact) \\
$2$ & $16$ & $14$ & $14$ & $0.0$ (exact) \\
$3$ & $40$ & $55$ & $55$ & $0.0$ (exact) \\
\end{tabular}
\end{center}
\medskip

The residual is exactly $0.0$ in float64 (not merely within machine
epsilon) because the two constructions call the same underlying
matrix-element routine and the $1 \times 1$ tensor with identity is
the identity bit-operation in IEEE 754.
\end{proof}

\begin{remark}[Significance of bit-exactness]\label{rem:bit_exact}
Bit-exact Riemannian-limit recovery is the natural consistency
condition for any Lorentzian extension that purports to reduce to its
Riemannian counterpart at static slices. Failure here would indicate
a structural incompatibility between the L2-B Krein-space basis
convention (Peskin--Schroeder chiral basis at $(m, n) = (4, 6)$) and
the FullDirac spinor-bundle basis convention (Paper~32~\S~III, Sprint
WH1-R3.5~\cite{paper32}). The passing of the falsifier confirms the
two conventions are bit-identical at $\Nt = 1$, and the Lorentzian
construction inherits the Riemannian spinor lift verbatim.
\end{remark}

% =====================================================================
\section{Envelope dependence of the propagation number}
\label{sec:propagation}
% =====================================================================

The propagation number~\cite{connes_vs2021}~Definition~2.39 of an
operator system $\Op \subset \Mat_N(\C)$ is the smallest $k \in \N
\cup \{\infty\}$ such that
$\dim_\C \Op^k = \dim_\C \mathrm{env}(\Op)$, where
$\Op^k = \mathrm{span}_\C \{a_1 \cdot a_2 \cdots a_k :\ a_i \in \Op\}$
and $\mathrm{env}(\Op)$ is the chosen envelope (typically the
$C^*$-envelope, equal to $\Mat_N(\C)$ in the Connes--vS Toeplitz
setting). On the chirality-doubled spinor bundle there are two
natural envelopes:\ the \emph{Weyl-doubled achievable} envelope and
the \emph{full} envelope. The principal result of this paper is
that the propagation number distinguishes them.

\subsection{The two envelopes}\label{sec:envelopes}

A scalar-multiplier product $M_1 \cdot M_2$ on $\HGV^{\nmax}$
preserves the chirality-block-diagonal structure:
\begin{equation}\label{eq:chirality_diag}
   (M_1 \oplus M_1) \cdot (M_2 \oplus M_2)
   \;=\; (M_1 M_2) \oplus (M_1 M_2).
\end{equation}
The maximal subspace of $\Bcal(\HGV^{\nmax})$ reachable by
scalar-multiplier products is therefore the Weyl-doubled subspace
\begin{equation}\label{eq:weyl_doubled}
   \mathcal{V}_\Weyl
   \;:=\; \bigl\{\, M_W \oplus M_W \;:\; M_W \in \Mat_{\dim_\Weyl}(\C)
                \,\bigr\}
   \;\cong\; \Mat_{\dim_\Weyl}(\C),
\end{equation}
with $\dim_\C \mathcal{V}_\Weyl = \dim_\Weyl^2$. Tensored with the
commutative temporal subalgebra
$\mathcal{D}_{\Nt} \subset \Mat_{\Nt}(\C)$ of dimension $\Nt$, the
achievable subspace on $\Krein_{\nmax, \Nt}$ is
\begin{equation}\label{eq:achievable_def}
   \mathcal{V}_\achievable
   \;:=\; \mathcal{V}_\Weyl \otimes \mathcal{D}_{\Nt}
   \;\subset\; \Bcal(\Krein),
\end{equation}
of dimension
\begin{equation}\label{eq:dim_achievable}
   \dim_\C \mathcal{V}_\achievable
   \;=\; \dim_\Weyl^2 \cdot \Nt.
\end{equation}
By contrast, the \emph{full} $\Bcal(\Krein)$ envelope has
\begin{equation}\label{eq:dim_full}
   \dim_\C \Bcal(\Krein_{\nmax, \Nt})
   \;=\; \dim_\Krein^2
   \;=\; \dim_\spat^2 \cdot \Nt^2
   \;=\; 4 \dim_\Weyl^2 \cdot \Nt^2,
\end{equation}
where the factor of $4$ comes from $\dim_\spat = 2 \dim_\Weyl$ on the
chirality-doubled bundle.

\begin{definition}[Envelope-relative propagation number]
\label{def:prop_env}
For envelopes $\mathcal{V} \in
\{\mathcal{V}_\achievable, \Bcal(\Krein)\}$, the propagation number
$\prop_{\mathcal{V}}(\Op^L)$ is the smallest $k \in \N \cup \{\infty\}$
such that $\dim_\C \mathrm{span}_\C(\Op^L)^k = \dim_\C \mathcal{V}$,
where $(\Op^L)^k = \mathrm{span}_\C\{a_1 a_2 \cdots a_k\}$ is the
$k$-fold product space. We say $\prop = \infty$ when no finite $k$
reaches $\dim_\C \mathcal{V}$.
\end{definition}

\subsection{Achievable-envelope propagation: prop = 2}
\label{sec:prop_ach}

\begin{theorem}[Achievable-envelope propagation]
\label{thm:prop_envelope_dependent}
For every $(\nmax, \Nt)$ with $\nmax \ge 2$ and $\Nt \ge 1$:
\begin{equation}\label{eq:prop_ach}
   \prop_{\achievable}(\Op_{\nmax, \Nt}^L) \;=\; 2.
\end{equation}
For every $(\nmax, \Nt)$ with $\nmax \ge 2$ and $\Nt \ge 1$ at finite
cutoff:
\begin{equation}\label{eq:prop_full_inf}
   \prop_{\mathrm{full}}(\Op_{\nmax, \Nt}^L) \;=\; \infty.
\end{equation}
At $\nmax = 1$, the spatial slot is trivial (only the
identity multiplier survives), and $\prop_{\achievable} = \infty$ as
well, the trivial case.
\end{theorem}

\begin{proof}[Proof of \eqref{eq:prop_full_inf}]
Every product $M_1 \cdot M_2$ of two scalar-multiplier generators is
chirality-block-diagonal by \eqref{eq:chirality_diag}, so
$(\Op^L)^k \subset \mathcal{V}_\Weyl \otimes
\mathrm{span}_\C\{M_p \cdot M_{p'} \cdots\}$ for every
$k$. Since $\mathcal{V}_\Weyl \subsetneq \Mat_{\dim_\spat}(\C)$ (the
off-block-diagonal chirality-flipping operators are not in
$\mathcal{V}_\Weyl$), and the temporal $C^*$-algebra generated by
diagonal matrices is $\Mat_{\Nt}(\C) \cap \mathcal{D}_{\Nt} =
\mathcal{D}_{\Nt}$ (the commutative subalgebra is its own $C^*$-%
algebra image), we have $(\Op^L)^k \subset \mathcal{V}_\achievable$
for every $k$. Therefore
$\dim_\C (\Op^L)^k \le \dim_\C \mathcal{V}_\achievable =
\dim_\Weyl^2 \cdot \Nt$, strictly less than
$\dim_\C \Bcal(\Krein) = \dim_\Krein^2 = 4 \dim_\Weyl^2 \cdot \Nt^2$
for any $\Nt \ge 1$ and $\nmax \ge 2$. So
$\prop_{\mathrm{full}}(\Op^L) = \infty$.
\end{proof}

\begin{proof}[Proof of \eqref{eq:prop_ach}, computational verification]
The Riemannian Paper~32~\S~III result is $\prop(\Op_{\nmax}^\Weyl) =
2$ for the Weyl-sector operator system on $\HGV^{\Weyl, \nmax}$
(Proposition~\ref{prop:paper32_prop}, see~\S~\ref{sec:context}). Under
the chirality-doubling map $\Op_{\nmax}^\Weyl \to
\Op_{\nmax}^{\mathrm{Full\,Dirac}}$, $M_W \mapsto M_W \oplus M_W$,
the propagation number is preserved with target envelope
$\mathcal{V}_\Weyl$ (the doubled image space):\
$(\Op^{\mathrm{Full\,Dirac}})^2$ fills $\mathcal{V}_\Weyl$ because
$(\Op^\Weyl)^2$ fills $\Mat_{\dim_\Weyl}(\C) =
\mathcal{V}_\Weyl$. Tensoring with the commutative
$\mathcal{D}_{\Nt}$:\ products of generators
$(M^\spat_a \otimes M_p^\temp) \cdot (M^\spat_b \otimes M_{p'}^\temp)
= (M^\spat_a M^\spat_b) \otimes (M_p^\temp M_{p'}^\temp)$ generate
$(M^\spat_a M^\spat_b)$-content in $\mathcal{V}_\Weyl$ tensored with
products $M_p^\temp M_{p'}^\temp \in \mathcal{D}_{\Nt}$. Since
$\{M_p^\temp\}_{p = 0}^{\Nt - 1}$ already spans $\mathcal{D}_{\Nt}$,
$(\Op^L)^2$ fills $\mathcal{V}_\Weyl \otimes \mathcal{D}_{\Nt} =
\mathcal{V}_\achievable$.

The computational verification is documented in
\texttt{tests/test\_operator\_system\_lorentzian.py::%
test\_propagation\_number\_nmax\_2\_*\_achievable\_envelope} and
\texttt{::test\_propagation\_number\_nmax\_3\_achievable\_envelope}.
A complete numerical panel is:

\medskip
\begin{center}\renewcommand{\arraystretch}{1.20}
\begin{tabular}{c|c|c|c|c|c|c}
$\nmax$ & $\Nt$ & $\dim_\Krein$ & $\dim_\Weyl$
   & $\dim_\achievable$
   & $\dim(\Op^L)$ & $\dim((\Op^L)^2)$ \\
\hline
$1$ & $1$ & $4$  & $2$ & $4$  & $1$  & $1$  \\
$1$ & $3$ & $12$ & $2$ & $12$ & $3$  & $3$  \\
$2$ & $1$ & $16$ & $8$ & $64$ & $14$ & $\mathbf{64}$ \\
$2$ & $3$ & $48$ & $8$ & $192$ & $42$ & $\mathbf{192}$ \\
$2$ & $5$ & $80$ & $8$ & $320$ & $70$ & $\mathbf{320}$ \\
$3$ & $1$ & $40$ & $20$ & $400$ & $55$ & $\mathbf{400}$ \\
\end{tabular}
\end{center}
\medskip

The bolded entries show $\dim((\Op^L)^2) = \dim_\achievable$, i.e.\
$\prop = 2$ at every cell with $\nmax \ge 2$. The $\nmax = 1$ cells
are the trivial case (the spatial slot has only the identity
multiplier $(N, L, M) = (0, 0, 0)$, so the operator system is
$\C \cdot I_\spat \otimes \mathcal{D}_{\Nt}$ of dimension $\Nt$,
which is closed under multiplication and so has no nontrivial
propagation depth to drive).
\end{proof}

\begin{proposition}[Paper~32~\S~III propagation, Weyl sector]
\label{prop:paper32_prop}
For every $\nmax \in \{2, 3, 4\}$,
$\prop(\Op_{\nmax}^\Weyl) = 2$ on the envelope
$\mathrm{env}(\Op^\Weyl) = \Mat_{\dim_\Weyl}(\C)$, with
$\dim(\Op^\Weyl) < \dim_\Weyl^2 = \dim(\mathrm{env})$ strictly.
\end{proposition}

\begin{proof}
Paper~32~\S~III, \texttt{prop:propagation\_number}, computational
verification in \texttt{tests/test\_operator\_system.py} (24 tests
passing, $\sim 20$ s CPU time). The result is invariant under three
SO(4) radial-overlap placeholder conventions (Paper~32~\S~III,
discussion after the proposition statement).
\end{proof}

\subsection{Why both numbers are real}\label{sec:both_real}

The discrepancy between $\prop_\achievable = 2$ and
$\prop_{\mathrm{full}} = \infty$ is not a technical glitch:\ it reflects
the fact that scalar multipliers on a chirality-doubled bundle
genuinely cannot reach the full $\Bcal(\Krein)$ envelope at any finite
$k$. The two natural targets correspond to two different operational
questions:

\medskip\noindent
\textbf{Achievable.} At what depth $k$ does the operator system fill
the maximal subspace structurally consistent with the
scalar-multiplier construction (chirality-block-diagonal, temporal-%
commutative)?  Answer: $k = 2$.

\medskip\noindent
\textbf{Full.} At what depth $k$ does the operator system reach the
full matrix algebra $\Mat_{\dim_\Krein}(\C)$?  Answer: never.
Reaching this envelope would require either Cl$(3, 1)$-gamma-matrix-%
valued multipliers (chirality-flipping) or a non-commutative temporal
multiplier algebra (e.g.\ $\partial_t$-extended), both of which are
outside the natural scalar-multiplier construction.

The Riemannian Paper~32~\S~III situation has no analogous distinction:\
on the scalar Fock-projected $\sthree$ Hilbert space (Paper~32
Definition~3.2), there is no chirality doubling and no temporal
factor, so the achievable and full envelopes coincide. The
distinction is genuinely introduced by the Lorentzian extension:\ the
chirality doubling enters at the spinor level (Sprint WH1-R3.5
\texttt{full\_dirac\_operator\_system.py}, which the Riemannian
Paper~32~\S~III analysis chose not to invoke for clarity); the
temporal factor enters at the L2-B Krein-space level. Both
contribute structural blocks to the propagation that are visible in
the present construction but invisible in the Riemannian template.

\begin{remark}[Connes--vS dual operator system as comparison]
\label{rem:comparison}
Connes--vS~\cite{connes_vs2021}~Proposition~4.3 establishes that the
\emph{dual} operator system to the Toeplitz $C(S^1)^{(n)}$, namely
$C(\Z)^{(n)} \subset \Mat_n(\C)$ generated by characteristic
functions on the integer lattice, has $\prop = \infty$. Our
$\prop_{\mathrm{full}} = \infty$ result has a different mechanism
(envelope-side, not generator-side), but the qualitative
phenomenon — propagation number depends on which natural envelope is
targeted — has a published parallel.
\end{remark}

\begin{remark}[Connes--vS \S IV sharpening:\ envelope-relative
informative content]\label{rem:cvs_sharpening}
The Connes--van~Suijlekom~\cite{connes_vs2021} reading of the
Riemannian truncated operator system, as articulated in
Paper~32~\S~III \texttt{rem:operator\_system}~\cite{paper32}, frames
the $C^*$-envelope $C^*(\Op_\Lambda)$ as the ``non-informative full
matrix algebra'' that the spectral truncation collapses to, while
the operator-system data $\Op_\Lambda$ itself carries the surviving
algebraic content. On the Toeplitz worked example
$C(S^1)^{(n)}$~\cite[\S\,3.1]{connes_vs2021} and on the
$\sthree$~scalar Fock-projected Hilbert space (Paper~32 Definition~3.2),
these two readings cohere into a single statement:\
$\dim(\Op_{\nmax}) < \dim(\mathrm{env}) = N^2$ strictly (the operator
system is informative because it is sub-envelope), yet
$\dim(\Op_{\nmax}^2) = N^2$ (its second power fills the envelope), so
$\prop = 2$. The Paper~32 witness pair (raise-then-lower at the top
shell) and the propagation-2 result are then two facets of the same
single envelope.

The present Lorentzian construction sharpens this Riemannian
template into a dichotomy that has no analog in
Paper~32~\S~III. Under the Weyl-doubled achievable envelope
$\mathcal{V}_\achievable$ — the structural target for
scalar-multiplier closure on the chirality-doubled spinor bundle —
the Paper~32 reading transports verbatim:\ the witness pair of
Theorem~\ref{thm:witness_pair} exhibits closure failure at
$\sim 38\%$ relative Frobenius residual at $\nmax = 2$ ($35.89\%$ at
$\nmax = 3$, \texttt{debug/data/l3a\_1\_lorentzian\_operator\_system.json}),
yet $\dim((\Op^L)^2) = \dim_\achievable$ exactly
(Theorem~\ref{thm:prop_envelope_dependent} eq.~\eqref{eq:prop_ach}).
The operator system is ``informative'' in the Connes--vS sense
because its products fill the natural target, and the witness pair
is the residual visible only inside that target.

Under the full $\Bcal(\Krein)$ envelope, however, the witness-pair
residual is supplemented by a deeper, second structural obstruction:\
chirality-block-diagonal closure plus a commutative temporal
subalgebra together cannot generate the chirality-flipping or
non-diagonal-temporal directions of $\Bcal(\Krein)$ at any finite
$k$ (Theorem~\ref{thm:prop_envelope_dependent}
eq.~\eqref{eq:prop_full_inf}). This is a real structural ceiling
imposed by the Lorentzian geometry itself — the chirality doubling
forced by the BBB $(m, n) = (4, 6)$ Krein structure plus the
commutative temporal slot — not a residual closure failure of the
truncation.

The operator system's informative content is therefore
\emph{envelope-relative}. On the achievable envelope, Paper~44
preserves Paper~32's Connes--vS structural reading exactly:\
sub-envelope linear dimension, full second-power dimension, $\prop =
2$. On the full envelope, the operator system is ``even more
sub-informative'' than its Riemannian template, because it hits the
Lorentzian-geometric ceiling before reaching $\Bcal(\Krein)$. We
therefore name the achievable envelope as the natural target for the
Connes--vS propagation question on the Lorentzian extension, with
the full-envelope $\prop = \infty$ recorded as a structural ceiling
that does not vitiate the achievable-envelope reading.
\end{remark}

% =====================================================================
\section{Trivial Krein-positive restriction}\label{sec:krein_positivity}
% =====================================================================

The Krein-positive cone is the $(+1)$-eigenspace of $\JL$:
\begin{equation}\label{eq:krein_plus}
   \Krein^+ \;:=\; \{ \psi \in \Krein \;:\; \JL \psi = +\psi \},
   \qquad \dim_\C \Krein^+ = \tfrac{1}{2} \dim_\C \Krein.
\end{equation}
An operator $A \in \Bcal(\Krein)$ \emph{preserves} $\Krein^+$ iff
$A \cdot \Krein^+ \subset \Krein^+$, equivalently iff
$[\JL, A] = 0$ at the level of the block decomposition
$\Krein = \Krein^+ \oplus \Krein^-$. The Krein-positive
restriction of $\Op^L$ is
\begin{equation}\label{eq:O_L_plus}
   \Op^{L, +} \;:=\; \{ A \in \Op^L \;:\; [\JL, A] = 0 \}.
\end{equation}

\begin{proposition}[Triviality of the Krein-positive restriction]
\label{prop:krein_positivity_trivial}
At every $(\nmax, \Nt)$ with $\Nt \ge 1$,
\begin{equation}\label{eq:O_L_plus_trivial}
   \Op^{L, +}_{\nmax, \Nt} \;=\; \Op_{\nmax, \Nt}^L.
\end{equation}
That is, every scalar-multiplier generator in $\Op^L$ already
preserves the Krein-positive cone.
\end{proposition}

\begin{proof}
A generator $M^{\spat}_{NLM} \otimes M_p^{\temp}$ commutes with
$\JL = \gamma^0 \otimes I_{\Nt}$ iff
$M^{\spat}_{NLM}$ commutes with $\gamma^0$ and
$M_p^{\temp}$ commutes with $I_{\Nt}$, the latter being automatic.

For the spatial factor, $M^{\spat}_{NLM} = M_W \oplus M_W$ in the
chirality decomposition $\HGV = \HGV^{\Weyl} \oplus \HGV^{\Weyl}$.
The fundamental symmetry $\gamma^0$ acts as the chirality-swap
matrix
\[
   \gamma^0
   \;=\; \begin{pmatrix} 0 & I_{\dim_\Weyl} \\
                          I_{\dim_\Weyl} & 0 \end{pmatrix}
\]
in the same decomposition (Paper~43~\S~3.1, Peskin--Schroeder chiral
basis). Direct computation:
\[
   \gamma^0 \cdot (M_W \oplus M_W)
   \;=\; \begin{pmatrix} 0 & I \\ I & 0 \end{pmatrix}
         \begin{pmatrix} M_W & 0 \\ 0 & M_W \end{pmatrix}
   \;=\; \begin{pmatrix} 0 & M_W \\ M_W & 0 \end{pmatrix},
\]
and
\[
   (M_W \oplus M_W) \cdot \gamma^0
   \;=\; \begin{pmatrix} M_W & 0 \\ 0 & M_W \end{pmatrix}
         \begin{pmatrix} 0 & I \\ I & 0 \end{pmatrix}
   \;=\; \begin{pmatrix} 0 & M_W \\ M_W & 0 \end{pmatrix},
\]
so $[\gamma^0, M^{\spat}_{NLM}] = 0$ exactly. Tensoring with any
operator commuting with $I_{\Nt}$ (everything),
$[\JL, M^{\spat}_{NLM} \otimes M_p^{\temp}] = 0$. The
computational verification is documented in
\texttt{tests/test\_operator\_system\_lorentzian.py::%
test\_krein\_positive\_restriction\_trivial}:\ at every tested
$(\nmax, \Nt) \in \{(2, 1), (2, 3), (3, 1)\}$, the fraction of
Krein-positive-preserving generators is exactly $1.0$.
\end{proof}

\begin{remark}[The Krein-positivity question shifts to the state level]
\label{rem:state_level}
Proposition~\ref{prop:krein_positivity_trivial} says that on the
chirality-doubled scalar-multiplier operator system, Krein
positivity is automatic at the level of the operators themselves.
The non-trivial Krein-positivity question is therefore not at the
operator-substrate level but at the \emph{state} level:\ which states
on $\Op^L$ assign positive expectation values to the Krein-positive
matrix elements?  The Wasserstein--Kantorovich state space on
$\Krein^+$ is the natural setting (cf.\
\cite{latremoliere_metric_st_2017}~\S~3, the Riemannian state-space
propinquity), and a Lorentzian-propinquity convergence argument
(Sprint L3b, see~\S~\ref{sec:open}) would presumably work at the
state level rather than the operator-system substrate.

A non-trivial Krein-positive restriction at the substrate level would
require enlarging $\Op^L$ to include chirality-mixing operators
(gamma-matrix-valued multipliers) and/or non-commutative temporal
operators (e.g., the Lorentzian Dirac $\DL$ itself, which is not in
$\Op^L$ as constructed). Both are structurally distinct constructions
from the natural pointwise-multiplication operator system studied
here.
\end{remark}

% =====================================================================
\section{The witness pair}\label{sec:witness}
% =====================================================================

\begin{theorem}[Witness pair for non-multiplicative-closure]
\label{thm:witness_pair}
Fix $\nmax \ge 2$, $\Nt \ge 1$. Let
\begin{equation}\label{eq:witness_a}
   a \;:=\; M^{\spat}_{N=2,\, L=1,\, M=0} \otimes M_0^{\temp}
   \;\in\; \Op_{\nmax, \Nt}^L,
\end{equation}
and $b := a^\dagger$. Then $a, b \in \Op_{\nmax, \Nt}^L$ but
$ab \notin \Op_{\nmax, \Nt}^L$:\ the least-squares projection
residual
\[
   \frac{\Frob{ab - \mathrm{proj}_{\Op^L}(ab)}}{\Frob{ab}}
\]
is bounded below by a positive constant uniformly in $\Nt$. The
numerical values at the tested cells are:

\medskip
\begin{center}\renewcommand{\arraystretch}{1.20}
\begin{tabular}{c|c|c}
$\nmax$ & $\Nt$ & relative Frobenius residual \\
\hline
$2$ & $1$ & $0.3812$ \\
$2$ & $3$ & $0.3812$ \\
$2$ & $5$ & $0.3812$ \\
$3$ & $1$ & $0.3589$ \\
\end{tabular}
\end{center}
\medskip

\noindent
The witness pair certifies that $\Op^L$ is genuinely an operator
system and not a $*$-algebra:\ it is closed under conjugate
transposition and linear combinations but \emph{not} under
multiplication.
\end{theorem}

\begin{proof}
At $p = 0$, $M_0^{\temp} = I_{\Nt}$, so
$a = M^{\spat}_{2,1,0} \otimes I_{\Nt}$ and
$ab = a a^\dagger = \bigl(M^{\spat}_{2,1,0} (M^{\spat}_{2,1,0})^\dagger\bigr)
      \otimes I_{\Nt}$.

The spatial factor
$M^{\spat}_{2,1,0} (M^{\spat}_{2,1,0})^\dagger$ is the Riemannian
witness product of Paper~32~\S~III~\cite{paper32}: raise-then-lower
at the top shell $n = \nmax$ produces a diagonal correction on the
top-shell entries that no scalar multiplier $M^{\spat}_{N L M}$ can
supply, because the corresponding $\nmax + 1$ shell is outside the
truncation~$P_{\nmax}$. Paper~32~\S~III reports a Frobenius residual
of $\sim 14.9\%$ at $\nmax = 2$ for the Weyl-sector witness; the
chirality-doubled FullDirac version inherits the same residual
structure with an arithmetic increase due to the doubled spatial
dimension, yielding the $0.3812$ value at $\nmax = 2$ recorded in
\texttt{debug/data/l3a\_1\_lorentzian\_operator\_system.json}. The
residual is independent of $\Nt$ at this label because $I_{\Nt}^2 =
I_{\Nt}$ is itself in $\mathcal{D}_{\Nt}$ (the constant temporal
multiplier), so the residual is purely spatial.

At $\nmax = 3$, the same construction with the same spatial label
$(N, L, M) = (2, 1, 0)$ produces a residual of $0.3589$:\ the
percentage drops slightly because the truncation reaches one shell
higher, partially absorbing the raise-then-lower correction. This
is the expected $\nmax$-scaling of the boundary deficit. The
verification is recorded in
\texttt{tests/test\_operator\_system\_lorentzian.py::%
test\_witness\_pair\_nmax\_2\_Nt\_1},
\texttt{::test\_witness\_pair\_nmax\_2\_Nt\_3}, and (for $\nmax = 3$)
\texttt{debug/data/l3a\_1\_lorentzian\_operator\_system.json}.
\end{proof}

\begin{remark}[Memo--JSON discrepancy at $\nmax = 3$]
\label{rem:nmax3_residual}
The Sprint L3a-1 memo~\S~5 quotes the witness-pair residual as
$0.381$ ``identical across $\Nt$ values'' across a panel that
includes $\nmax = 3, \Nt = 1$. The supporting JSON data
\texttt{debug/data/l3a\_1\_lorentzian\_operator\_system.json},
panel cell $(\nmax, \Nt) = (3, 1)$, records $0.3589$ instead. The
discrepancy is purely a memo bookkeeping issue;\ the JSON and tests
agree that the $\nmax = 2$ cells give $0.3812$ and the $\nmax = 3$
cell gives $0.3589$, with the slightly lower number reflecting the
expected boundary-deficit attenuation at the larger Fock cutoff. We
report both the $\nmax = 2$ ($0.3812$) and $\nmax = 3$ ($0.3589$)
values explicitly in Theorem~\ref{thm:witness_pair} per the JSON,
which is the data of record.
\end{remark}

% =====================================================================
\section{Connection to the Paper~43 hemispheric wedge}\label{sec:wedge}
% =====================================================================

Paper~43~\cite{paper43} works with the Krein hemispheric wedge
\begin{equation}\label{eq:W_L_def}
   W_L \;:=\; P_W^{\spat} \otimes P_{t \ge 0}
   \;\subset\; \Bcal(\Krein_{\nmax, \Nt}),
\end{equation}
where $P_W^{\spat}$ is the Paper~42~\cite{paper42} hemispheric
Hopf-axis $m_j$-reflection projector on $\HGV^{\nmax}$ and
$P_{t \ge 0}$ is the diagonal projector on $\C^{\Nt}$ selecting
$t_k \ge 0$. The wedge-restricted operator system
\begin{equation}\label{eq:O_L_wedge}
   \Op_{\nmax, \Nt}^{L, W} \;:=\; \{\, P_{W_L} A P_{W_L}
                                     \;:\; A \in \Op^L \,\}
\end{equation}
is a sub-operator-system on the wedge sub-Krein-space
$\Krein_{W_L} := P_{W_L} \Krein \cap P_{W_L} \Krein$ of dimension
\begin{equation}\label{eq:dim_wedge}
   \dim_\C \Krein_{W_L}
   \;=\; \dim_\C \Krein_W^{\spat} \cdot \dim_\C \mathrm{range}(P_{t \ge 0})
   \;=\; \tfrac{1}{2} \dim_\spat \cdot \lceil \Nt / 2 \rceil
\end{equation}
(Paper~43~\S~4.1; for the symmetric grid, $\dim_\C \mathrm{range}(P_{t
\ge 0}) = \lceil \Nt / 2 \rceil$).

\begin{proposition}[Wedge-restricted operator-system data]
\label{prop:wedge_restriction}
At every $(\nmax, \Nt)$ with $\nmax \ge 2, \Nt \ge 1$, the
wedge-restricted operator system
$\Op_{\nmax, \Nt}^{L, W}$ is a well-defined sub-operator-system on
$\Krein_{W_L}$. The wedge linear-span dimensions at the tested cells
are:

\medskip
\begin{center}\renewcommand{\arraystretch}{1.20}
\begin{tabular}{c|c|c|c|c}
$\nmax$ & $\Nt$ & $\dim_\C \Krein_{W_L}$
   & $\dim_\C \Op^{L,W}$ & comment \\
\hline
$2$ & $1$ & $8$  & $10$  & wedge linear span exceeds wedge dim \\
$2$ & $3$ & $16$ & $20$  & same \\
$2$ & $5$ & $24$ & $30$  & same \\
\end{tabular}
\end{center}
\medskip

The wedge linear-span dimension exceeds the wedge Hilbert-space
dimension because $\Op^{L,W}$ is a subspace of
$\Bcal(\Krein_{W_L}) \cong \Mat_{\dim_{W_L}}(\C)$ whose dimension is
$\dim_{W_L}^2$ ($64$ at the $\nmax = 2, \Nt = 1$ cell), not
$\dim_{W_L}$.
\end{proposition}

\begin{proof}
The wedge projection is a direct application of the
\texttt{LorentzianWedge} construction from
\texttt{geovac/modular\_hamiltonian\_lorentzian.py} (Paper~43~\S~4.1).
The wedge-restricted matrix is computed by index-restriction:
$(P_{W_L} A P_{W_L})|_{i, j \in W_L} = A_{i, j}$ for indices $i, j$
in the wedge basis-index set. The linear-span dimension is computed by
SVD on the stack of vec-form wedge-restricted multiplier matrices,
documented in
\texttt{tests/test\_operator\_system\_lorentzian.py::%
test\_restrict\_to\_lorentzian\_wedge\_*}.
\end{proof}

\begin{remark}[Paper~43 modular content reads on the wedge-restricted
operator system]
\label{rem:wedge_paper43}
Paper~43~\cite{paper43} constructs the wedge KMS state $\rho_W^L$
and the modular Hamiltonians $K_L^\alpha$ and $K_L^{\TT}$ at the
matrix level on $\Krein_{W_L}$. The natural operator-system reading
of those constructions is at the level of expectation values on
$\Op^{L, W}_{\nmax, \Nt}$:\ the cross-witness collapse of
Paper~43~\S~6 (Cor.~8.1) is the statement that the trace functional
$A \mapsto \Tr_{\Krein_{W_L}}(\rho_W^L \cdot A)$ on $\Op^{L, W}$ is
identical across all four physical witness instantiations. The
present paper's contribution is to identify the carrier of this
construction as the wedge restriction of the natural Lorentzian
truncated operator system, completing the bottom-up picture.

This identification does \emph{not} immediately imply that the
Paper~43 modular flow extends to a flow on the operator system
itself (in the sense of an automorphism group of $\Op^L$), because
$\Op^L$ is not multiplicatively closed and ``automorphism'' does not
literally apply. The natural Lorentzian-propinquity-style replacement
— a continuous family of $\C$-linear maps preserving the
operator-system data — is the natural next direction (Sprint L3b).
\end{remark}

% =====================================================================
\section{Open questions and Sprint L3b outlook}\label{sec:open}
% =====================================================================

This paper closes the operator-system substrate at finite cutoff. The
following questions remain open.

\subsection{Continuum Lorentzian propinquity convergence}
\label{sec:open_propinquity}

\textbf{Question.} Does the family
$\{\Tcal^L_{\nmax, \Nt}\}_{\nmax, \Nt \to \infty}$ converge in some
Lorentzian-propinquity-style topology to a continuum limit triple
$\Tcal^L_\infty$ on $\sthree \times \R$?

This is the principal open question. The Riemannian
template — Paper~38~\cite{paper38}'s state-space Gromov--Hausdorff
convergence on $\sthree = \SU(2)$, Paper~39~\cite{paper39}'s
tensor-product extension, Paper~40~\cite{paper40_unified}'s
universal $4/\pi$ rate constant on compact Lie groups — relies on
five lemmas (L1$'$ chirality-doubled operator system, L2 central
spectral Fejér kernel, L3 Lipschitz spinor bound, L4 Berezin
reconstruction, L5 Latrémolière propinquity assembly). The
$\sthree \times \R$ analog requires the construction of a joint
compact $\times$ non-compact Fejér kernel and a Krein-positive
propinquity definition, both of which are mathematically open
(\texttt{debug/l3\_scoping\_memo.md}~\S~4, Q1 and Q4). As of May 2026,
no published Lorentzian propinquity has been constructed:\
Latr\'emoli\`ere's 2018~\cite{latremoliere2018} and
2017/2023~\cite{latremoliere_metric_st_2017} works are strictly
Riemannian, as are
Hekkelman--McDonald~\cite{hekkelman_mcdonald2024,
hekkelman_mcdonald2024b} and
Farsi--Latr\'emoli\`ere~\cite{farsi_latremoliere2024,
farsi_latremoliere2025}.

\medskip\noindent
\textbf{Recommended first move (Sprint L3b).} Replace the non-compact
$\R_t$ by a compact $S^1_T$ with circumference $T < \infty$,
producing a joint compact $\SU(2) \times S^1_T$ group manifold on
which Peter--Weyl analysis applies in both factors. The L2 central
spectral Fejér kernel then transports via a product construction with
an explicit joint rate. This reduces L3b to a tractable 4--8 week
sprint (\texttt{debug/l3\_scoping\_memo.md}~\S~6). The
compactified-temporal target is then the natural setting for a
Lorentzian-propinquity assembly that inherits all five Paper~38 lemmas
factor-by-factor, modulo Krein-positivity conventions at the
Berezin reconstruction step.

\subsection{Krein-positivity at the state level}\label{sec:open_state}

\textbf{Question.} What is the natural Wasserstein--Kantorovich
state-space distance on the Krein-positive states of $\Op^L$, and
does it converge under the Sprint L3b construction?

Proposition~\ref{prop:krein_positivity_trivial} shifts the
Krein-positivity question from the operator-substrate to the state
level. The natural object is the set of positive linear functionals
$\omega:\ \Op^L \to \C$ with $\omega(A^\dagger A) \ge 0$ for $A \in
\Op^L$ \emph{and} $\omega$ respects the Krein indefinite form (i.e.\
$\omega(\JL A \JL) = \omega(A)$, the natural Krein invariance). The
Wasserstein--Kantorovich distance on this state space is then the
candidate Lorentzian-propinquity metric. To our knowledge, no
construction of such an object appears in the published NCG literature
as of May 2026.

\subsection{Nieuviarts twist-morphism applicability}\label{sec:open_nieuviarts}

\textbf{Question.} Does the Nieuviarts~2025~\cite{nieuviarts2025a}
twist morphism from Riemannian-twisted to pseudo-Riemannian spectral
triples apply to $\sthree = \SU(2)$ at KO-dimension $3$?

Nieuviarts~\cite{nieuviarts2025a}~Definition~2.2 explicitly restricts
the twist morphism to even-dimensional manifolds. The author's
companion proceedings note~\cite{nieuviarts2025b_proceedings}
acknowledges that the odd-dimensional case is open. The present
construction explicitly bypasses this by working directly with the
van~den~Dungen~2016~\cite{vandungen2016} Proposition~4.1 lift to a
Krein space, which has no parity-dimension restriction. If the
Nieuviarts construction were extended to odd dimensions, it would
provide an alternative substrate for the Sprint L3b program;\ as it
stands, the BBB Krein-construction path of the present paper is the
only path available.

\subsection{Multi-focal-composition wall}\label{sec:open_multifocal}

\textbf{Question.} Does the framework's structural-skeleton scope
limitation (Paper~42~\S~10 O5; CLAUDE.md~\S~1.7 multi-focal-wall
taxonomy; \texttt{debug/multifocal\_phase\_b\_synthesis\_memo.md},
2026-05-08) apply to the Lorentzian extension?

This is open and structurally independent of the present paper;\ the
W2a renormalization wall, the W2b cross-manifold wall, and the W3
calibration-data question (the ``second packing axiom''
problem) are not addressed by Sprint L3 or the present paper. They
remain in the same status as the Riemannian closure of Paper~42.

\subsection{Higher-rank compact Lie group extension}\label{sec:open_higher_rank}

\textbf{Question.} For compact $G$ in the Paper~40~\cite{paper40_unified}
class ($G \in \{\SU(N), \mathrm{Spin}(n), \mathrm{Sp}(n), G_2, F_4,
E_6, E_7, E_8\}$), does the Lorentzian extension $G \times \R$ admit
the same operator-system substrate with envelope-dependent
propagation?

The construction in the present paper uses only that
$\HGV^{\nmax}$ is a chirality-doubled spinor bundle on a compact
group manifold with multiplier action via chirality-block doubling;
the same construction transports immediately to higher-rank $G$ with
the spinor bundle replaced by the appropriate $G$-spinor module.
The propagation-number invariant
$\prop_\achievable = 2$ should hold at every $G$ by the same
chirality-block argument, since the underlying Paper~32~\S~III
prop$ = 2$ result extends to higher-rank $G$ via
Paper~40~\cite{paper40_unified}'s L1$'$ lemma. Verification at
$\SU(3)$, $\SU(4)$, $\mathrm{Sp}(2)$, $G_2$ would be a natural
$\sim 4$ week extension of the present module.

\subsection{Cross-manifold W2b}\label{sec:open_w2b}

\textbf{Question.} Does the Lorentzian extension transport to the
cross-manifold target $\Tcal_{\sthree}^L \otimes
\Tcal_{\mathrm{Hardy}(S^5)}^L$?

The cross-manifold tensor product is the W2b open question of the
multi-focal-composition wall taxonomy (CLAUDE.md~\S~1.7), structurally
blocked at the NCG-framework level by the Paper~24~\cite{paper24}~\S~V
Coulomb/HO category mismatch:\ the Camporesi--Higuchi spinor bundle
on $\sthree$ is a second-order Riemannian construction, while the
Bargmann--Segal Hardy-space construction on $S^5$ is a first-order
complex-analytic construction. The category mismatch resurfaces at
the spectral-triple level (Paper~32~\S~VIII.C). The present paper
does not address W2b;\ it remains open.

% =====================================================================
\section{Conclusion}\label{sec:conclusion}
% =====================================================================

We have generalized the Connes--van~Suijlekom~\cite{connes_vs2021}
truncated operator system from the Riemannian round $\sthree =
\SU(2)$ Camporesi--Higuchi spectral triple~\cite{camporesi_higuchi1996}
(Paper~32~\S~III~\cite{paper32}) to the Bizi--Brouder--Besnard~%
\cite{bizi_brouder_besnard2018} Krein spectral triple at signature
$(s, t) = (3, 1)$ in the Peskin--Schroeder chiral basis at
$(m, n) = (4, 6)$, with Lorentzian Dirac
$\DL = i(\gamma^0 \otimes \partial_t + \DGV \otimes I_{\Nt})$ per
van~den~Dungen~2016~\cite{vandungen2016} Proposition~4.1. The
construction reduces bit-identically to the chirality-doubled
Riemannian operator system on $\HGV^{\nmax}$ at $\Nt = 1$, confirming
compatibility with the Sprint L2-B Krein-space basis convention
(Theorem~\ref{thm:riemannian_limit}).

The principal structural finding is the
\emph{envelope dependence of the propagation number}
(Theorem~\ref{thm:prop_envelope_dependent}). Targeting the
Weyl-doubled achievable envelope $\dim_\Weyl^2 \cdot \Nt$, the
Lorentzian extension preserves the Connes--vS propagation-2 result
verbatim (eq.~\eqref{eq:prop_ach}, matching Paper~32~\S~III, matching
the Toeplitz operator system $C(S^1)^{(n)}$ of
\cite{connes_vs2021}~Proposition~4.2). Targeting the full
$\Bcal(\Krein)$ envelope, the propagation number is $\infty$
(eq.~\eqref{eq:prop_full_inf}):\ chirality-doubling and a
commutative temporal subalgebra block the scalar-multiplier
construction from generating chirality-flipping operators or
non-diagonal temporal operators. This envelope-dependence is a real
structural feature of the chirality-doubled Lorentzian extension,
not visible in the scalar Riemannian template.

A separate substantive finding is that the Krein-positive restriction
to the natural sub-operator-system $\Op^{L, +}$ equals $\Op^L$
itself (Proposition~\ref{prop:krein_positivity_trivial}). Every
chirality-doubled scalar multiplier $M \oplus M$ commutes with the
chirality-swap fundamental symmetry $\JL$, so Krein positivity is
automatic at the operator-substrate level. The non-trivial
Krein-positivity question is at the \emph{state} level, where the
Wasserstein--Kantorovich state-space construction on Krein-positive
states is the natural setting for a future Lorentzian-propinquity
argument.

\medskip\noindent
\textbf{Honest scope.} This paper closes the operator-system
substrate at finite cutoff. It does \emph{not} establish:
\begin{enumerate}\setlength{\itemsep}{2pt}
\item Continuum / GH-limit / Lorentzian-propinquity convergence
(open; Sprint L3b is the recommended first move,
\S~\ref{sec:open_propinquity}).
\item Krein-positive state-space convergence
(open; \S~\ref{sec:open_state}).
\item Non-commutative temporal multipliers (the natural pointwise-%
multiplication algebra on $\R_t$ is commutative; this is the
canonical Connes--vS construction).
\item Gamma-matrix-valued multipliers (which would enlarge the
operator system from $C^\infty(\sthree \times \R)$ to
$C^\infty(\sthree \times \R, \Cl(\R^4))$ and give a non-trivial
Krein-positive restriction).
\end{enumerate}

\medskip\noindent
\textbf{Position in the GeoVac Lorentzian arc.} Together with
Paper~43~\cite{paper43} (the four-witness Wick-rotation theorem at
the Lorentzian Krein-wedge level) and the prospective Sprint L3b
continuum closure, the present paper constitutes the operator-system
substrate of the GeoVac Lorentzian extension at the level the
framework can currently reach. Paper~43 does not supersede the
present paper, nor vice versa:\ Paper~43 establishes the
modular-Hamiltonian content on a hemispheric wedge; the present
paper establishes the operator-system content on the full Krein
space and its envelope-relative structural invariants. Both are
finite-cutoff closures with the multi-month continuum problem named
explicitly and left open.

\medskip\noindent
\textbf{Position against the broader literature.} The Connes--van~%
Suijlekom~\cite{connes_vs2021} program deferred the GH-convergence
question on a non-trivial example to ``elsewhere'' three times.
Paper~38~\cite{paper38} closed it on the Riemannian
$\sthree = \SU(2)$ at the qualitative-rate state-space
Gromov--Hausdorff level; Paper~40~\cite{paper40_unified} closed it on the universal
compact-Lie-group family with rate constant $4/\pi$. The Lorentzian
analog of Paper~38 (continuum Lorentzian propinquity on
$\sthree \times \R$) remains open. The present paper, together with
Paper~43~\cite{paper43}, establishes the finite-cutoff substrate from
which any such convergence argument would proceed.

% =====================================================================
\section*{Acknowledgments}
% =====================================================================
The author thanks Alain Connes, Walter van~Suijlekom, Fr\'ed\'eric
Latr\'emoli\`ere, Nadir Bizi, Christian Brouder, Fabien Besnard,
Koen van~den~Dungen, Arnaud Nieuviarts, and the SageMath, sympy,
mpmath, and numpy communities for the mathematical and computational
infrastructure on which this paper builds; remaining errors are the
author's own. This work was developed in an AI-augmented agentic
workflow with Anthropic's Claude:\ implementation of the supporting
computational module
\texttt{geovac/operator\_system\_lorentzian.py}, drafting of the
internal proof memo
\texttt{debug/l3a\_1\_lorentzian\_operator\_system\_memo.md}, and
bibliographic collation were performed collaboratively. All
theorem statements, design choices, and editorial decisions are the
author's responsibility.

% =====================================================================
\bibliographystyle{plainnat}
\begin{thebibliography}{99}

\bibitem[Wen \& Avery(1985)]{avery_wen_avery2002}
S.~Wen and J.~Avery,
\newblock \emph{Some products of spherical harmonics and related
integrals},
\newblock J.~Math.~Phys. \textbf{26} (1985), 396--403.

\bibitem[Bisognano \& Wichmann(1976)]{bisognano_wichmann1976}
J.~J.~Bisognano and E.~H.~Wichmann,
\newblock \emph{On the duality condition for quantum fields},
\newblock J.~Math.~Phys. \textbf{17} (1976), 303--321.

\bibitem[Bizi, Brouder \& Besnard(2018)]{bizi_brouder_besnard2018}
N.~Bizi, C.~Brouder, and F.~Besnard,
\newblock \emph{Space and time dimensions of algebras with
application to Lorentzian noncommutative geometry and quantum
electrodynamics},
\newblock J.~Math.~Phys. \textbf{59} (2018), 062303;
\newblock arXiv:1611.07062.

\bibitem[Bykov, Minguzzi \& Suhr(2024)]{bykov_minguzzi_suhr2024}
A.~Bykov, E.~Minguzzi, and S.~Suhr,
\newblock \emph{Lorentzian metric spaces and GH-type convergence},
\newblock arXiv:2412.04311, 2024.

\bibitem[Camporesi \& Higuchi(1996)]{camporesi_higuchi1996}
R.~Camporesi and A.~Higuchi,
\newblock \emph{On the eigenfunctions of the Dirac operator on
spheres and real hyperbolic spaces},
\newblock J.~Geom.~Phys. \textbf{20} (1996), 1--18.

\bibitem[Connes \& van~Suijlekom(2021)]{connes_vs2021}
A.~Connes and W.~D.~van~Suijlekom,
\newblock \emph{Spectral truncations in noncommutative geometry and
operator systems},
\newblock Comm.~Math.~Phys. \textbf{383} (2021), 2021--2067;
\newblock arXiv:2004.14115.

\bibitem[Devastato, Lizzi \& Martinetti(2018)]{devastato_lizzi_martinetti2018}
A.~Devastato, F.~Lizzi, and P.~Martinetti,
\newblock \emph{Lorentzian almost-commutative spectral triples},
\newblock J.~Math.~Phys. \textbf{59} (2018), 092304.

\bibitem[Farsi \& Latr\'emoli\`ere(2024)]{farsi_latremoliere2024}
C.~Farsi and F.~Latr\'emoli\`ere,
\newblock \emph{Collapse in noncommutative geometry and spectral
continuity},
\newblock arXiv:2404.00240, 2024.

\bibitem[Farsi \& Latr\'emoli\`ere(2025)]{farsi_latremoliere2025}
C.~Farsi and F.~Latr\'emoli\`ere,
\newblock \emph{Continuity for the spectral propinquity of the
Dirac operators associated with an analytic path of Riemannian
metrics},
\newblock arXiv:2504.11715, 2025.

\bibitem[Franco \& Eckstein(2014)]{franco_eckstein2014}
N.~Franco and M.~Eckstein,
\newblock \emph{An algebraic formulation of causality for
noncommutative geometry},
\newblock Class.~Quantum Grav. \textbf{30} (2013), 135007;
\newblock arXiv:1212.5171.

\bibitem[Hartle \& Hawking(1976)]{hartle_hawking1976}
J.~B.~Hartle and S.~W.~Hawking,
\newblock \emph{Path-integral derivation of black-hole radiance},
\newblock Phys.~Rev. D \textbf{13} (1976), 2188--2203.

\bibitem[Hekkelman(2022)]{hekkelman2022}
E.-M.~Hekkelman,
\newblock \emph{Truncated Geometry on the Circle},
\newblock arXiv:2111.13865, 2022.

\bibitem[Hekkelman \& McDonald(2024)]{hekkelman_mcdonald2024}
E.~Hekkelman and E.~McDonald,
\newblock \emph{Spectral truncations of $T^d$ and quantum metric
geometry},
\newblock J.~Noncommut.~Geom., to appear (2024);
\newblock arXiv:2403.18619.

\bibitem[Hekkelman \& McDonald(2024b)]{hekkelman_mcdonald2024b}
E.~Hekkelman and E.~McDonald,
\newblock \emph{A noncommutative integral on spectrally truncated
spectral triples, and a link with quantum ergodicity},
\newblock J.~Funct.~Anal., to appear (2025);
\newblock arXiv:2412.00628.

\bibitem[Latr\'emoli\`ere(2017)]{latremoliere_metric_st_2017}
F.~Latr\'emoli\`ere,
\newblock \emph{The Gromov--Hausdorff propinquity for metric spectral
triples},
\newblock Adv.~Math. \textbf{404} (2022), Paper No.~108393, 56pp;
\newblock preprint arXiv:1811.10843, 2017.

\bibitem[Latr\'emoli\`ere(2018)]{latremoliere2018}
F.~Latr\'emoli\`ere,
\newblock \emph{The Gromov--Hausdorff propinquity},
\newblock Trans.~Amer.~Math.~Soc. \textbf{370} (2018), 365--411.

\bibitem[Leimbach \& van~Suijlekom(2024)]{leimbach_vs2024}
M.~Leimbach and W.~D.~van~Suijlekom,
\newblock \emph{Gromov--Hausdorff convergence of spectral truncations
of the torus},
\newblock Adv.~Math. \textbf{439} (2024), Paper No.~109496.

\bibitem[Mondino \& Sämann(2025)]{mondino_samann2025}
A.~Mondino and C.~Sämann,
\newblock \emph{Lorentzian Gromov--Hausdorff convergence and
pre-compactness},
\newblock arXiv:2504.10380, 2025.

\bibitem[Nieuviarts(2025a)]{nieuviarts2025a}
G.~Nieuviarts,
\newblock \emph{Emergence of Lorentz symmetry from an
almost-commutative twisted spectral triple},
\newblock arXiv:2502.18105 (v3, May 2025).

\bibitem[Nieuviarts(2025b)]{nieuviarts2025b_proceedings}
G.~Nieuviarts,
\newblock \emph{Emergence of time from a twisted spectral triple
in almost-commutative geometry},
\newblock arXiv:2512.15450, December 2025 (revised May 2026).

\bibitem[Paulsen(2002)]{paulsen2002}
V.~I.~Paulsen,
\newblock \emph{Completely Bounded Maps and Operator Algebras},
\newblock Cambridge Studies in Advanced Mathematics 78,
Cambridge University Press, 2002.

\bibitem[Sewell(1982)]{sewell1982}
G.~L.~Sewell,
\newblock \emph{Quantum fields on manifolds:\ PCT and gravitationally
induced thermal states},
\newblock Ann.~Phys. (NY) \textbf{141} (1982), 201--224.

\bibitem[Strohmaier(2006)]{strohmaier2006}
A.~Strohmaier,
\newblock \emph{On noncommutative and pseudo-Riemannian geometry},
\newblock J.~Geom.~Phys. \textbf{56} (2006), 175--195.

\bibitem[Toyota(2023)]{toyota2023}
M.~Toyota,
\newblock \emph{Quantum Gromov--Hausdorff convergence of spectral
truncations for groups with polynomial growth},
\newblock arXiv:2309.13469, 2023.

\bibitem[Unruh(1976)]{unruh1976}
W.~G.~Unruh,
\newblock \emph{Notes on black-hole evaporation},
\newblock Phys.~Rev. D \textbf{14} (1976), 870--892.

\bibitem[van~den~Dungen(2016)]{vandungen2016}
K.~van~den~Dungen,
\newblock \emph{Krein spectral triples and the fermionic action},
\newblock Math.~Phys.~Anal.~Geom. \textbf{19} (2016), Art.~4, 22pp;
\newblock arXiv:1505.01939.

\bibitem[Loutey, internal Paper~24]{paper24}
J.~Loutey,
\newblock \emph{The Bargmann--Segal Lattice: $\pi$-Free Discretization
of the Harmonic Oscillator on $S^5$},
\newblock GeoVac internal preprint (2026), DOI via Zenodo.

\bibitem[Loutey, internal Paper~32]{paper32}
J.~Loutey,
\newblock \emph{The GeoVac Spectral Triple: A Synthesis Paper Locking
the Framework into the Marcolli--van~Suijlekom Gauge-Network Lineage},
\newblock GeoVac internal preprint, papers/synthesis (2026),
DOI via Zenodo.

\bibitem[Loutey, internal Paper~38]{paper38}
J.~Loutey,
\newblock \emph{State-space Gromov--Hausdorff convergence of
truncated Camporesi--Higuchi spectral triples on $S^3$},
\newblock GeoVac internal preprint, papers/standalone (2026),
DOI via Zenodo.

\bibitem[Loutey, internal Paper~39]{paper39}
J.~Loutey,
\newblock \emph{Tensor-product propinquity convergence for
distinct-focal-length truncated Camporesi--Higuchi spectral triples
on $\sthree$},
\newblock GeoVac internal preprint, papers/standalone (2026),
DOI via Zenodo.

\bibitem[Loutey, internal Paper~40]{paper40_unified}
J.~Loutey,
\newblock \emph{State-space Gromov--Hausdorff convergence of spectral
truncations of compact Lie groups with bi-invariant metric},
\newblock GeoVac internal preprint, papers/standalone (2026),
DOI via Zenodo.

\bibitem[Loutey, internal Paper~42]{paper42}
J.~Loutey,
\newblock \emph{Tomita--Takesaki modular structure on truncated
$SU(2)$ spectral triples: four-witness Wick-rotation literal
identification at finite cutoff},
\newblock GeoVac internal preprint, papers/standalone (2026),
DOI via Zenodo.

\bibitem[Loutey, internal Paper~43]{paper43}
J.~Loutey,
\newblock \emph{Lorentzian extension of the four-witness
Wick-rotation theorem on truncated $S^3 \times \mathbb{R}$ spectral
triples at finite cutoff},
\newblock GeoVac internal preprint, papers/standalone (2026),
DOI via Zenodo.

\end{thebibliography}

% =====================================================================
\section*{Appendix: Section reference for the \texorpdfstring{$\S$~\ref{sec:propagation}}{Sec. 5} context cross-reference}
\label{sec:context}
% =====================================================================

The Riemannian Paper~32~\S~III propagation-number result
(Proposition~\ref{prop:paper32_prop}) was established by direct
computation in \texttt{geovac/operator\_system.py} (the
\texttt{TruncatedOperatorSystem} class). The verification covers
$\nmax \in \{2, 3, 4\}$ with multiplier counts
$\dim_\C \Op_{\nmax} = 14, 55, 140$ and propagation-table

\medskip
\begin{center}\renewcommand{\arraystretch}{1.20}
\begin{tabular}{c|c|c|c|c|c}
$\nmax$ & $N_{\mathrm{Fock}}$ & $N^2$ & $\dim(\Op)$ & $\dim(\Op^2)$
   & $\dim(\Op) / N^2$ \\
\hline
$2$ & $5$  & $25$  & $14$  & $25$  & $0.560$ \\
$3$ & $14$ & $196$ & $55$  & $196$ & $0.281$ \\
$4$ & $30$ & $900$ & $140$ & $900$ & $0.156$ \\
\end{tabular}
\end{center}
\medskip

reproduced from Paper~32~\S~III, \texttt{prop:propagation\_number}.
The result is robust under three SO(4) radial-overlap convention
choices (including the indicator-only convention, which retains only
the support pattern of the matrix-element bookkeeping; see
\texttt{debug/wh1\_round2\_propagation\_memo.md}, 2026-05-03). The
present paper's Theorem~\ref{thm:prop_envelope_dependent} inherits
this robustness verbatim:\ the achievable-envelope propagation number
$\prop_\achievable = 2$ is the same SO(4) selection-rule structural
invariant lifted to the chirality-doubled Lorentzian construction.

\end{document}
